\section{Experiments}
\label{sec:experiments}

\subsection{Dataset and Evaluation Metrics}
\label{sec:dataset}

The CGPS dataset~\cite{cv4chl2026} contains 2D keypoint annotations from multi-view clinical video recordings of 110 pediatric patients at Shenzhen Children's Hospital. Each frame provides 133 COCO-WholeBody~\cite{jin2020whole} keypoints detected by Sapiens-2B~\cite{khirodkar2024sapiens} with bounding boxes produced using the challenge organizers' segmentation pipeline~\cite{cv4chl2026}. The dataset spans 339,236 frames across 1,185 videos captured from four viewpoints (left sagittal, right sagittal, forward coronal, backward coronal) at 1920$\times$1080 resolution. Track~1 includes 78 training and 16 test patients; Track~2 includes 18 training and 9 test patients.

\noindent\textbf{Evaluation Metrics.}
For Track~1, the score combines item-level accuracy and normalized RMSE of total scores:
$S_1 = (\text{Acc}_1 + 1 - \text{RMSE}/34) / 2$.
For Track~2, the score combines limb-level accuracy and macro F1:
$S_2 = (\text{Acc}_2 + F1_{\text{macro}}) / 2$.
The originally announced composite score was $S = (S_1 + S_2) / 2$. The organizers subsequently reported a post-challenge weighting of $S = 0.7S_1 + 0.3S_2$ for their final re-evaluation, owing to Track~2's small-sample bias~\cite{cv4chl2026}. Accordingly, our recorded version scores below are historical competition feedback and must not be interpreted as scores from the organizers' final standardized re-evaluation.


\subsection{Evaluation Protocol and Interpretation}
\label{sec:protocol}
The CGPS recordings represent multiple views and potentially multiple trials of the same patient; hence, frame-level or video-level random splitting would risk patient-identity leakage. The appropriate independent unit for model selection is the patient. We report the originally recorded solution-version feedback as descriptive evidence, not as a new patient-disjoint controlled benchmark. The official organizers later performed three standardized evaluations, including alternative train/test partitions, and revised the relative track weights~\cite{cv4chl2026}. This difference in evaluation protocol limits comparisons between our historical version scores and the final official rankings.

\noindent\textbf{Data scope and label semantics.}
The source dataset contains 110 pediatric participants, 1,185 video recordings and 339,236 annotated frames, with 133 WholeBody landmarks per frame. The challenge supplies pose sequences and metadata instead of releasing the original patient videos. Track~1 binarizes the three-level clinical EVGS categories by mapping moderate and severe deviations to the positive class; Track~2 predicts one of five gait patterns for each limb. Our reported training/test counts refer to the challenge task subsets, not the full set of 110 patients.

\noindent\textbf{Reproducibility boundary.}
The solution-version comparisons were conducted during iterative competition development. Because some variants jointly changed preprocessing, features, model settings and decision rules, the reported differences should be treated as observational comparisons rather than independent causal estimates. No confidence intervals, external hospital validation or prospective clinical assessment are available in the present study.

\subsection{Implementation Details}
\label{sec:implementation}

All models are implemented using scikit-learn~\cite{pedregosa2011scikit}. We use shallow tree-based ensembles: Gradient Boosting (100 estimators, max depth~3, learning rate~0.1), Random Forest (200 estimators, max depth~6), and Extra Trees (200 estimators, max depth~6). Feature selection retains the top-50 features for Track~1 and top-26 for Track~2, determined by stabilized importance across 20 seeds. The hybrid weight $\alpha = 0.6$ is selected through leaderboard validation. Leave-One-Out Cross-Validation (LOO-CV) is used for internal evaluation, though we note its unreliability for this dataset due to distribution shift (\cf Section~\ref{sec:laws}).

\subsection{Main Results}
\label{sec:main_results}

Table~\ref{tab:main_results} presents the leaderboard scores across our solution versions, ordered chronologically. Our best submission (V7) achieves a combined score of 0.6664.

\begin{table}[t]
    \centering
    \small
    \caption{\textbf{Leaderboard scores across solution versions.} Each version introduces specific methodological changes. The best result (V7) combines domain knowledge integration, aggressive cleaning, and cross-version ensemble.}
    \label{tab:main_results}
    \begin{tabular}{clc}
        \toprule
        Ver. & Key Strategy & Score \\
        \midrule
        V1 & Stat.\ features + GB/RF/ET baseline & 0.498 \\
        V2 & + Clinical angles, sagittal view & 0.499 \\
        V4 & + Patient-level T2, stable importance & 0.590 \\
        V5 & + Data cleaning (single smooth) & 0.625 \\
        V6 & V5 cleaning + V4 structure (combo) & 0.576 \\
        \textbf{V7} & \textbf{+ Rule hybrid, double smooth, ens.} & \textbf{0.666} \\
        V9 & + Gaussian-calibrated rules & 0.647 \\
        V10 & + 100-seed ensemble & 0.556 \\
        V11 & KNN zero-training baseline & 0.574 \\
        V12c & Pure rule-based T2 & 0.416 \\
        V13 & Item-specific T1 features & 0.645 \\
        \bottomrule
    \end{tabular}
\end{table}

\subsection{Track Decomposition Analysis}
\label{sec:decomposition}

To understand the contribution of each track, we conducted isolation experiments by combining one track from V7 with the corresponding track from V1 (baseline). Table~\ref{tab:decomposition} shows that 97\% of our total improvement originates from Track~2 ($\Delta S_2 = 0.326$), while Track~1 barely improved ($\Delta S_1 = 0.011$). This reveals that Track~1 remains the primary bottleneck.

\begin{table}[t]
    \centering
    \small
    \caption{\textbf{Track decomposition via isolation experiments.} V7's improvement over baseline is almost entirely from Track~2.}
    \label{tab:decomposition}
    \begin{tabular}{lcc}
        \toprule
        Configuration & LB Score & $\Delta$ vs.\ V1 \\
        \midrule
        V1 (baseline) & 0.498 & --- \\
        V7-T1 + V1-T2 & 0.503 & $\Delta S_1 \approx +0.011$ \\
        V1-T1 + V7-T2 & 0.661 & $\Delta S_2 \approx +0.326$ \\
        V7 (both tracks) & 0.666 & $+0.168$ \\
        \bottomrule
    \end{tabular}
\end{table}

\subsection{Ablation Studies}
\label{sec:ablation}

We systematically ablate each component of our framework. Table~\ref{tab:ablation} summarizes the contribution of each key design choice.

\begin{table}[t]
    \centering
    \small
    \caption{\textbf{Ablation study.} Incremental contribution of each component, estimated from controlled version comparisons.}
    \label{tab:ablation}
    \begin{tabular}{lc}
        \toprule
        Component & Contribution \\
        \midrule
        Baseline (V1) & 0.498 \\
        + Patient-level T2 classification & +0.092 \\
        + Single-pass temporal cleaning & +0.035 \\
        + Double-pass temporal cleaning & +0.041 \\
        + 40\% rule hybrid for T2 & +0.076 \\
        + Cross-track EVGS features & +0.035 \\
        + Cross-version ensemble & +0.030 \\
        \midrule
        Best (V7) & 0.666 \\
        \bottomrule
    \end{tabular}
\end{table}

\noindent\textbf{Data Cleaning.}
Data cleaning is the single most impactful factor. Without cleaning, scores plateau around 0.50--0.59. Single-pass smoothing (V5) improves to 0.625, and double-pass smoothing (V7) further improves to 0.666. Crucially, cleaning shows zero or negative impact on LOO-CV but massive impact on the leaderboard, confirming that the test set is significantly noisier.

\noindent\textbf{Rule Integration Layer.}
The layer at which domain knowledge is integrated matters critically. Rules as features (V5) add $+0.035$; rules at the prediction layer with 40\% weight (V7) add $+0.076$. However, pure rules without ML (V12c) score only 0.416---worse than the baseline---because angles in normalized coordinate space do not directly match clinical thresholds. The optimal operating point lies at approximately 40\% rule weight.

\noindent\textbf{Ensemble Granularity.}
Moderate ensembling (7 seeds $\times$ 3 models $=$ 21 classifiers) is effective, but excessive ensembling (100 seeds $\times$ 3 models $=$ 300 classifiers in V10) \textit{degrades} performance by 0.110, as over-averaging drives probabilities toward uniform distributions and destroys discriminative power.

\subsection{Empirical Laws for Small-Sample Clinical ML}
\label{sec:laws}

Through systematic experimentation across 14 solution versions, we identify 12 empirical laws. Table~\ref{tab:laws} provides a summary. We highlight three key findings below.

\begin{table}[t]
    \centering
    \small
    \caption{\textbf{Twelve empirical laws} derived from systematic experimentation. Confidence rated on a 5-star scale based on the number of supporting experiments.}
    \label{tab:laws}
    \resizebox{\linewidth}{!}{
    \begin{tabular}{clc}
        \toprule
        \# & Empirical Law & Conf. \\
        \midrule
        1 & Data cleaning is the \#1 LB driver & 5/5 \\
        2 & Hardcoded rules $>$ data-fitted rules & 4/5 \\
        3 & Rules at prediction layer $>$ rules as features & 5/5 \\
        4 & Naive combination of good methods can hurt & 5/5 \\
        5 & LOO-CV is anti-correlated with LB score & 5/5 \\
        6 & Patient-level T2 classification is effective & 4/5 \\
        7 & Moderate ensemble $>$ excessive ensemble & 5/5 \\
        8 & KNN fails for Track~1 (78 patients) & 3/5 \\
        9 & 97\% of improvement comes from Track~2 & 5/5 \\
        10 & Item-specific T1 feature restriction hurts & 3/5 \\
        11 & Pure rules fail in normalized coordinate space & 4/5 \\
        12 & Cross-version ensemble adds genuine diversity & 3/5 \\
        \bottomrule
    \end{tabular}
    }
\end{table}

\noindent\textbf{Law 5: Cross-Validation is Anti-Correlated with Test Performance.}
The most counter-intuitive finding is that LOO-CV scores are negatively correlated with leaderboard performance. V4 achieves CV$=$0.681 but LB$=$0.590, while V5 achieves CV$=$0.487 but LB$=$0.625. This arises because training and test sets have different noise distributions---methods robust to noise score poorly on the cleaner training set but generalize better to the noisier test set. This fundamentally challenges the use of cross-validation for model selection in small clinical datasets with distribution shift.

\noindent\textbf{Law 4: Naive Combination of Good Methods Can Hurt.}
Combining V5's data cleaning with V4's feature selection (V6) yields 0.576---\textit{worse than either method alone} (V5$=$0.625, V4$=$0.590). Similarly, combining item-specific features with V7's Track~2 strategy (V13) yields 0.645, below V7's 0.666. Each method's internal assumptions are optimized for its own feature space; merging them creates conflicts that degrade performance.

\noindent\textbf{Law 3: Rules Integrated at the Prediction Layer Outperform Rules as Features.}
Using clinical rules as input features (V5, $+0.035$) is less effective than hybrid fusion at the prediction layer (V7, $+0.076$). When rules are encoded as features, the ML model can ``dilute'' or override the clinical signal through overfitting on the small training set. Prediction-layer integration preserves clinical knowledge as an independent signal source that cannot be overridden.

\subsection{Gap Analysis}
\label{sec:gap}

Our best recorded score (0.666) should not be directly compared against the organizers' later post-challenge ranking, which used additional evaluation splits and a revised track weighting. Our within-submission decomposition suggests that most of our measured gain over the V1 baseline arose from Track~2, leaving Track~1 as a key opportunity for further improvement. We hypothesize that top-performing teams likely employ: (1)~deep learning architectures (\eg, ST-GCN~\cite{yan2018spatial} or Transformers) directly on raw keypoint sequences, enabling end-to-end temporal feature learning; (2)~pre-training on external gait datasets; (3)~multi-task learning jointly optimizing Track~1 and Track~2; and (4)~sophisticated data augmentation including temporal warping and viewpoint synthesis. Our handcrafted feature approach fundamentally limits the model's capacity to discover discriminative patterns beyond predefined joint angles.
